% concatenated from main.tex, macros.tex
\documentclass{article}
\usepackage{graphicx} % Required for inserting images
\usepackage{amsthm}
\usepackage{amsfonts, amsmath, amssymb, amsthm}
\usepackage{hyperref}
\usepackage{cleveref}
\usepackage[utf8x]{inputenc}
\usepackage{url}
\usepackage{tikz-cd}

\def\TT{\mathbb{T}}
\def\PP{\mathbb{P}}
\def\DD{\mathbb{D}}
\def\QQ{\mathbb{Q}}
\def\WW{\mathbb{W}}
\def\MM{\mathbb{M}}
\def\HH{\mathbb{H}}

\newcommand{\Psf}{\mathsf{P}}
\newcommand{\Qsf}{\mathsf{Q}}
\newcommand{\Rsf}{\mathsf{R}}


\def\F{\mathcal{F}}
\def\U{\mathcal{U}}
\def\C{\mathcal{C}}
\def\I{\mathcal{I}}
\def\D{\mathcal{D}}
\def\L{\mathcal{L}}
\def\V{\mathcal{V}}
\def\A{\mathcal{A}}
\def\B{\mathcal{B}}
\def\P{\mathcal{P}}
\def\M{\mathcal{M}}
\def\Nc{\mathcal{N}}
\def\Q{\mathcal{Q}}

% Baire space, Cantor space, binary strings
\newcommand{\cs}{2^{\NN}}
\newcommand{\bs}{\NN^{\NN}}
\newcommand{\bstr}{2^{<\NN}}

\newcommand{\uh}{\upharpoonright}

% Turing machines
\newcommand{\converge}{\!\!\downarrow}
\newcommand{\diverge}{\!\!\uparrow}

% Forcing question
\newcommand{\qvdash}{\operatorname{{?}{\vdash}}}
\newcommand{\nqvdash}{\operatorname{{?}{\nvdash}}}

\newcommand{\ISig}{\mathsf{I}\Sigma^0}
\newcommand{\Robinson}{\mathsf{Q}}
\newcommand{\IPi}{\mathsf{I}\Pi^0}
\newcommand{\IDelta}{\mathsf{I}\Delta^0}
\newcommand{\Iopen}{\mathsf{I_{open}}}
\newcommand{\BSig}{\mathsf{B}\Sigma^0}
\newcommand{\BPi}{\mathsf{B}\Pi^0}
\newcommand{\BDelta}{\mathsf{B}\Delta^0}
\newcommand{\BCSig}{\mathsf{BC}\Sigma^0}
\newcommand{\BCPi}{\mathsf{BC}\Pi^0}
\newcommand{\BCDelta}{\mathsf{BC}\Delta^0}
\newcommand{\LSig}{\mathsf{L}\Sigma^0}
\newcommand{\LPi}{\mathsf{L}\Pi^0}
\newcommand{\CSig}{\mathsf{C}\Sigma}

\newcommand{\RCA}[0]{\mathsf{RCA}}
\newcommand{\WKL}[0]{\mathsf{WKL}}
\newcommand{\WWKL}[0]{\mathsf{WWKL}}
\newcommand{\ACA}[0]{\mathsf{ACA}}
\newcommand{\ATR}[0]{\mathsf{ATR}}
\newcommand{\PIOOCA}[0]{\Pi^1_1\mathsf{-CA}}
\newcommand{\CAC}[0]{\mathsf{CAC}}
\newcommand{\ADS}[0]{\mathsf{ADS}}
\newcommand{\RRT}[0]{\mathsf{RRT}}
\newcommand{\FS}[0]{\mathsf{FS}}
\newcommand{\PFS}[0]{\mathsf{PFS}}
\newcommand{\TS}[0]{\mathsf{TS}}
\newcommand{\GRRT}[0]{\mathsf{GRRT}}
\newcommand{\LRRT}[0]{\mathsf{LRRT}}
\newcommand{\EM}[0]{\mathsf{EM}}
\newcommand{\SEM}[0]{\mathsf{SEM}}
\newcommand{\HalfEM}[0]{\mathsf{Half}\mbox{-}\mathsf{EM}}
\newcommand{\RT}[0]{\mathsf{RT}}
\newcommand{\SRT}[0]{\mathsf{SRT}}
%\newcommand{\RRT}{\mathsf{RRT}}

\newcommand{\NN}[0]{\mathbb{N}}
\newcommand{\nat}[0]{\mathbb{N}}

\newcommand{\card}{\operatorname{card}}
\newcommand{\finsub}{\subseteq_{\mathtt{fin}}}

% reductions

\newcommand{\W}{\mathrm{W}}
\newcommand{\sW}{\mathrm{sW}}

\newcommand{\lex}{\mathrm{lex}}

% Barriers
\newcommand{\base}{\operatorname{base}}
\newcommand{\range}{\operatorname{ran}}

%typographically correct double quotes
\def\qt#1{``#1''}%

\title{Ramsey-like theorems for the Schreier barrier}
\date{\today}

\newtheorem*{statement}{Statement}
\newtheorem{theorem}{Theorem}
\numberwithin{theorem}{section}
\newtheorem{maintheorem}[theorem]{Main Theorem}
\newtheorem{lemma}[theorem]{Lemma}
\newtheorem{question}[theorem]{Question}
\newtheorem{proposition}[theorem]{Proposition}
%\newtheorem{statement}[theorem]{Statement}
\newtheorem{remark}[theorem]{Remark}
\newtheorem{definition}[theorem]{Definition}
\newtheorem{corollary}[theorem]{Corollary}
\newtheorem{example}[theorem]{Example}

\newcommand{\R}{\mathcal{R}}
\renewcommand{\S}{\mathcal{S}}
\newcommand{\T}{\mathcal{T}}


% Allow to repeat theorems
\makeatletter
\newtheorem*{rep@theorem}{\rep@title}
\newcommand{\newreptheorem}[2]{%
\newenvironment{rep#1}[1]{%
 \def\rep@title{#2 \ref{##1}}%
 \begin{rep@theorem}}%
 {\end{rep@theorem}}}
\makeatother

\newreptheorem{theorem}{Theorem}

% Comments
\usepackage{xcolor}
\newcommand{\ludovic}[1]{\textcolor{blue}{#1}}
\newcommand{\quentin}[1]{\textcolor{green}{#1}}
\newcommand{\lorenzo}[1]{\textcolor{orange}{#1}}
\newcommand{\oriola}[1]{\textcolor{red}{#1}}

\author{Lorenzo Carlucci \and Oriola Gjetaj \and Quentin Le Houérou \and Ludovic Levy Patey}

\title{Uniformly computable arrays of thin colourings whose thin sets compute the k-th Turing jump}
\author{Lorenzo Carlucci, Oriola Gjetaj, Quentin Le Houérou, Ludovic Levy Patey}
\date{arxiv:2412.11598v3; CC BY 4.0}
\begin{document}
\maketitle
\setcounter{section}{4}
\setcounter{theorem}{0}
\noindent\textit{Editorial note.} The counted unit is the article's own Lemma lem:arrays (source lines 683-685) together with the proof that begins on the next line (687-697); the pairing is the ordinary adjacent one.\par\medskip
\noindent\textit{Editorial note.} The sentence that motivates the lemma, at source line 679, contrasts this uniform construction with the non-uniform finite-dimensional result of Dorais et al.; it is not extracted because it carries a citation for which this fixed source supplies no bibliography entry, and its content is recorded here and in the allowed prerequisites instead.\par\medskip
\noindent\textit{Editorial note.} This article numbers all theorem-like environments in one sequence within each section, so the retained definitions print as 4.1 and 4.2 and the counted lemma as 4.3, matching the source's own numbering; the original label lem:arrays is carried in the span metadata.\par\medskip
\noindent\textit{Editorial note.} The authors' iffalse block at source lines 584-626 and the iffalse section at 1307-1376 are not extracted, so no counted unit rests on material the authors removed.\par\medskip
\noindent\textit{Editorial note.} No retained line contains a citation, a figure or an included file.\par\medskip
\noindent\textit{Editorial scope.} One original result: the article's Lemma lem:arrays with its complete author proof, plus the article's own definitions of an f-thin set and of the finite-dimensional Thin Set Theorem. The five claims already counted from this source, in particular the omega-th-jump theorem that later consumes this lemma, are not counted here and are not used as premises.\par\medskip
\noindent\textit{Allowed standard prerequisites (editorial).}\begin{itemize}
\item Turing functionals, jumps and moduli. The Turing functional Phi\_e\textasciicircum{}X, the k-th Turing jump of the empty set (the halting set for k = 1), the class Delta\textasciicircum{}0\_\{k+1\}, a uniform modulus of a set (a function dominating its characteristic function's use bound), and the fact that a function dominating a modulus of the k-th jump computes that jump. All of these are used as they stand at source lines 688 and 696 and are not re-proved.. The colouring is built from a uniform modulus and the conclusion is the computability of the k-th jump from a thin set.
\item Thin sets for a colouring. For a colouring f from the n-element subsets of the natural numbers into a non-empty set C and a colour c in C, an infinite H is f-thin for c when c is not in f applied to the n-element subsets of H; the article's own definition is retained from source lines 665-667, as is the principle TS\textasciicircum{}n\_k at 675-677.. The statement and every step of the proof are phrased as thinness for an explicitly named colour.
\item Finite combinatorics of tuples in an infinite set. An infinite set of natural numbers supplies arbitrarily long increasing tuples, and tuples can be completed beyond any prescribed bound while keeping the first coordinates fixed.. Used at source lines 692 and 694 to force the top colour and then to realise any smaller colour.
\end{itemize}\medskip
\begin{definition}
Let $C$ be any non-empty set. Given a coloring $f : [\NN]^n \to C$, a set~$H \subseteq \NN$ is \emph{$f$-thin for color~$c \in C$} if $c \not \in f([H]^n)$. A set ~$H \subseteq \NN$ is $f$-thin (or thin for $f$) if $H$ is $f$-thin for some color $c\in C$.
\end{definition}

\begin{definition}[Thin Set Theorem for finite colorings]\label{def:ts-finite}
Let $n, k\in \NN^+$ with $k\geq 2$. For every coloring $f : [\NN]^n \to k$, there exists an infinite set $H \subseteq \NN$ such that $H$ is thin for $f$. We abbreviate this statement by $\TS^n_{k}$.
\end{definition}

\begin{lemma}\label{lem:arrays}
    There exists two computable arrays $(e_{n,k})_{n,k \in \NN}$ and $(d_{n,k})_{n,k \in \NN}$ of Turing indexes such that for every $n, k \in \NN^+$ with $n \geq 2$, $\Phi_{e_{n,k}}^{\emptyset^{(k)}}$ is a coloring $f_n^k : [\NN]^{n} \to n$ such that for every infinite $f_n^k$-thin set $H$, $\Phi_{d_{n,k}}^{H} = \emptyset^{(k)}$.
\end{lemma}

\begin{proof}
Let $k\in\NN^+$. Let $g_k$ be a uniform modulus of the set $\emptyset^{(k)}$. Note that $\Delta^0_{k+1}$-indexes for each function $g_k$ can be found computably, uniformly in~$k$. 

For $n \geq 2$ let $f_n^k(x_0, \dots,x_{n-1}) = n - 1$ if $g_k(x_0) \leq x_1$ and if $g_k(x_0) > x_1$, let $f_n^k(x_0, \dots,x_{n-1}) = i$ for $i < n - 1$ the largest value such that $g_k(x_0) > x_{i+1}$.

Let $H$ be an infinite $f_n^k$-thin set for some color $c$. The color $c$ cannot be $n-1$ as for a given $x_0 \in H$ there exists $x_1 <\dots < x_{n-1} \in H \setminus \{0, \dots, g_k(x_0)\}$, hence $f^k_n(x_0, \dots, x_{n-1}) = n-1$.

We claim $H$ is thin for color $n-2$. Indeed, assume by contradiction that there exists some tuple $x_0 < \dots < x_{n-1} \in H$ such that $g_k(x_0) > x_{n-1}$ so that $f^k_n$ takes color $n-2$ on $\{x_0, \dots, x_{n-1}\}$. Since $c < n-2$ and $H$ is infinite, there exists some $y_{c+2}< \dots < y_{n-1} \in H \setminus \{0, \dots ,g_k(x_0)\}$, therefore $f^k_n(x_0, \dots, x_{c+1},y_{c+2}, \dots, y_{n-1}) = c$, contradicting our assumption that $H$ is $f_n^k$-thin for $c$.

Write $H = \{x_0 < x_1 < \dots \}$. For every $i \in \NN$, $g_k(x_i) \leq x_{i + n - 1}$, hence $H$ computes a function dominating $g_k$ and therefore computes $\emptyset^{(k)}$. This computation can be done uniformly in the set $H$, $k$ and $n$. 
\end{proof}

\end{document}
